% NNGraph DSL -- Persian LaTeX Report
% Compile: xelatex -shell-escape report-fa-humanized.tex  (run twice)

\documentclass{nngraph-fa}

\usepackage{xepersian}

\settextfont[
    BoldFont       = Vazirmatn Bold,
    ItalicFont     = Vazirmatn,
    BoldItalicFont = Vazirmatn Bold,
]{Vazirmatn}
\setdigitfont{Vazirmatn}
\setmonofont{Liberation Mono}

% ----------------------------------------------------------------
\title{
    \textbf{کامپایلر \LR{NNGraph DSL}} \\[0.5em]
    \large یک زبان دامنه‌محور برای توصیف معماری شبکه‌های عصبی \\
    و تولید خودکار کد \LR{PyTorch}
}
\author{\large گزارش فنی پروژه}
\date{\today}

% ----------------------------------------------------------------
\begin{document}

\maketitle
\tableofcontents
\newpage

% ================================================================
\chapter*{چکیده}
\markboth{چکیده}{چکیده}
\addcontentsline{toc}{chapter}{چکیده}

راستش \LR{NNGraph DSL} یه زبان دامنه‌محور باحاله که ما برای توصیف معماری شبکه‌های عصبی به شکل گراف جهت‌دار طراحیش کردیم
فایل‌های \code{.nng} رو می‌دیم بهش و اونم برامون کد اجرایی \LR{PyTorch} تولید می‌کنه که یه زیرکلاس \code{nn.Module} تر و تمیز با متدهای \code{\_\_init\_\_} و \code{forward} می‌ده بیرون

برای کامپایلر رفتیم سراغ \LR{ANTLR4} نسخه ۴.۱۳.۲ و از الگوی \LR{Listener} برای تحلیل معنایی استفاده کردیم
یه الگوریتم مرتب‌سازی توپولوژیک \LR{Kahn} هم زدیم تنگش که هم تو تحلیل‌گر معنایی (واسه تشخیص دور) و هم تو تولید کد (واسه ترتیب اجرا) حسابی به کارمون میاد
تازه استنتاج شکل تنسور، کشیدن گراف با \LR{Graphviz}، داربست آموزش، حاشیه‌نویسی کوانتیزاسیون و خروجی \LR{ONNX} رو هم بهش اضافه کردیم

تو بخش تست هم ۵۵ تا تست تو چهار تا فایل نوشتیم
چهار تا معماری باحال مثل \LR{MLP}، \LR{ResBlock}، \LR{InceptionBlock} و \LR{TransformerEncoder} رو هم کامپایل کردیم و \LR{forward pass} شون رو با شکل خروجی درست اجرا کردیم

% ================================================================
\chapter{مقدمه}

\section{انگیزه}

نوشتن زیرکلاس‌های \code{nn.Module} تو \LR{PyTorch} واسه معماری‌های پیچیده خیلی رو اعصابه و کلی باگ میده
مثلاً واسه یه شبکه ده لایه، ما باید:

\begin{itemize}
    \item تو \code{\_\_init\_\_} هر لایه رو دستی با پارامترای خودش تعریف کنیم
    \item تو \code{forward} هم ترتیب اجرا رو خودمون مدیریت کنیم
    \item تو توپولوژی‌های شاخه‌دار هم باید واسه متغیرهای میانی اسم بذاریم
    \item حواسمون باشه هیچ تنسور مشترکی قبل از اینکه استفاده بشه رو خودش بازنویسی نشه
\end{itemize}

یه نگاه به این کد \LR{PyTorch} بندازین که مجبوریم دستی بنویسیم:

\begin{codebox}{نمونه boilerplate دستی برای یک ResBlock}{code-resblock-manual}
\begin{amzcode}{python}
class ResBlock(nn.Module):
    def __init__(self):
        super(ResBlock, self).__init__()
        self.conv1 = nn.Conv2d(
            in_channels=64, out_channels=64,
            kernel_size=3, stride=1, padding=1)
        self.bn1   = nn.BatchNorm2d(num_features=64)
        self.relu1 = nn.ReLU()
        # ... and so on for each layer

    def forward(self, x):
        identity = x
        x = self.conv1(x)
        x = self.bn1(x)
        x = self.relu1(x)
        # ... manage residual connection manually
        x = x + identity
        return x
\end{amzcode}
\end{codebox}

با \LR{NNGraph DSL}، ما معماری رو یه بار مثل یه گراف توصیف می‌کنیم و کامپایلر بقیه کدها رو برامون تولید می‌کنه

\section{هدف}

هدف ما از \LR{NNGraph DSL} اینه که فایل \code{.nng} بشه تنها جایی که معماری رو توش نگه می‌داریم
اینطوری گراف رو توصیف می‌کنیم و کامپایلر کل اون کدهای تکراری رو خودش برامون می‌سازه

\section{دامنه کاربرد}

به نظرم این ابزار واسه معماری‌های ثابت خیلی جوابه، یعنی شبکه‌هایی که گراف محاسباتی‌شون همون اول کار معلومه
البته واسه نسخه‌های بعدی تو برنامه‌مون هست که کنترل جریان پویا تو \code{forward} رو هم اضافه کنیم

% ================================================================
\chapter{کارهای مرتبط}

\section{چارچوب NADER}

\begin{dfnbox}{NADER}{nader}
\dfntxt{\LR{NADER}} (\LR{neural architecture design via representation})
یه چارچوبه که ما باهاش معماری شبکه‌های عصبی رو به شکل یه گراف محاسباتی جهت‌دار (\LR{DAG}) مدل می‌کنیم
اینجا هر گره یه عملیات لایه‌ست و هر یال هم جریان تنسور رو نشون میده
\end{dfnbox}

ما تو \LR{NNGraph DSL} همین مدل گراف رو پایه کارمون قرار دادیم

\section{ANTLR4}

\begin{dfnbox}{ANTLR4}{antlr4}
\dfntxt{\LR{ANTLR4}} (\LR{another tool for language recognition})
یه ابزار خفنه که از رو فایل گرامر \code{.g4} برامون \LR{lexer}، \LR{parser} و \LR{listener/visitor} رو خودکار می‌سازه
نسخه‌ای که ما استفاده کردیم \LR{ANTLR 4.13.2} هست
\end{dfnbox}

\subsection{الگوی Listener در برابر Visitor}

کلاً \LR{ANTLR4} دو تا الگو واسه پیمایش درخت تجزیه داره:

\begin{itemize}
    \item \textbf{\LR{listener}}: اینجا \code{ParseTreeWalker} خودش درخت رو پیمایش می‌کنه و متدهای \code{enter*/exit*} رو صدا می‌زنه
    نیازی هم نیست چیزی برگردونه
    \item \textbf{\LR{visitor}}: تو این یکی ما خودمون پیمایش رو دست می‌گیریم و هر متدی یه چیزی برمی‌گردونه
\end{itemize}

\begin{notebox}
ما تو این کامپایلر از الگوی \LR{Listener} استفاده کردیم
متدهای \code{exit*} تو فایل \code{custom\_listener.py} وضعیت رو تو \code{self.nodes} و \code{self.edges} نگه می‌دارن، که به نظرم واسه ساختن گراف از رو درخت خیلی طبیعی‌تره
\end{notebox}

\section{ابزارها و تجسم‌سازهای AST}

واسه این کامپایلر ما از ابزارهای استاندارد درخت نحو انتزاعی (\LR{AST}) واسه بازرسی و دیدن درخت استفاده کردیم
کلاً سه تا فایل واسه این کار داریم:

\begin{itemize}
    \item \code{ast.py} --- کلاس \code{AST} و \code{TreeNode}
    \item \code{make\_ast\_subtree.py} --- ترکیب گره‌های درخت تجزیه
    \item \code{ast\_to\_networkx\_graph.py} --- کشیدن گراف به صورت اختیاری
\end{itemize}

\subsection{چرا AST در اینجا استفاده نمی‌شود}

بیشتر کامپایلرا برنامه‌ها رو به صورت درخت نحو انتزاعی (\LR{AST}) نشون می‌دن
زبان‌هایی که رو عبارت‌ها می‌چرخن، خیلی راحت به ساختار درختی تبدیل می‌شن

اما داستان \LR{NNGraph DSL} فرق داره
ما معماری شبکه‌های عصبی رو مثل یه گراف جهت‌دار بدون دور (\LR{DAG}) می‌بینیم
گره‌ها همون لایه‌ها هستن و یال‌ها هم جریان تنسورها رو نشون می‌دن
مدل ما اینجا خود گرافه، نه درخت تجزیه

واسه همین ما کلی عملیات گرافی انجام می‌دیم که به لیست مجاورت و یال‌ها نیاز دارن:

\begin{itemize}[noitemsep]
    \item \textbf{مرتب‌سازی توپولوژیک} (الگوریتم \LR{Kahn}) واسه ترتیب اجرا تو \code{forward()}
    \item \textbf{تشخیص دور} با مرتب‌سازی توپولوژیک (واسه اینکه مطمئن بشیم \LR{DAG} ما درسته)
    \item \textbf{تحلیل دسترس‌پذیری} با یه \LR{BFS} مشتی از ورودی تا خروجی
    \item \textbf{اعتبارسنجی اتصال باقی‌مانده} (چک کردن درجه ورودی گره‌های ادغام)
    \item \textbf{استنتاج شکل} با انتشار شکل تنسورها رو یال‌ها
    \item \textbf{خروجی \LR{Graphviz DOT}} واسه اینکه معماری رو قشنگ ببینیم
    \item \textbf{خروجی \LR{ONNX}} (که خودش بر اساس فرمت گرافه)
\end{itemize}

یه \LR{AST} فقط سلسله مراتب درخت تجزیه رو به ما می‌ده (\code{model\_block}، \code{graph\_block}، \code{node\_decl})، اما گراف محاسباتی که ما می‌خوایم رو قایم می‌کنه
واسه همین تحلیل‌گر معنایی ما گره‌ها و یال‌ها رو مستقیم تو دیکشنری‌ها و لیست‌ها ذخیره می‌کنه که واسه الگوریتم‌های گراف خیلی بهتره

\begin{notebox}
کدای ساخت \LR{AST} (\code{exitEveryRule}، \code{exitStart}) تو \code{custom\_listener.py} فقط واسه دیباگ اختیاری با پرچم \code{-{}-show-ast} اونجا هستن و موقع کامپایل اصلاً کاری به کارشون نداریم
\end{notebox}

% ================================================================
\chapter{طراحی زبان}

\section{ساختار برنامه}

هر فایل \code{.nng} کلاً سه تا بلوک داره
بلوک \code{config} هم که اختیاریه:

\begin{codebox}{ساختار کلی یک فایل .nng}{code-structure}
\begin{amzcode}{text}
model <Name> {
    input  <id> : tensor(<shape>)
    output <id>
}

graph {
    node <id> : <LayerType>(<params>)
    ...
    edge <src> -> <dst>
    edge <src> -> <dst> [label="<string>"]
    ...
}

config {
    batch_size = <int>
    device     = "cpu" | "cuda"
    loss       = "<LossFn>"       // enables training scaffold
    optimizer  = "<Optimizer>"    // Adam | SGD | AdamW | RMSprop
    lr         = <float>          // learning rate
    epochs     = <int>            // training epochs
}
\end{amzcode}
\end{codebox}

\begin{dfnbox}{بلوک model}{block-model}
بلوک \code{model} اسم کلاس تولیدشده، شناسه گره ورودی و شکل تنسور ورودی، و شناسه گره خروجی رو مشخص می‌کنه
\end{dfnbox}

\begin{dfnbox}{بلوک graph}{block-graph}
ما تو بلوک \code{graph} همه گره‌های لایه و یال‌های جهت‌دار رو تعریف می‌کنیم
ترتیب نوشتن یال‌ها اصلاً مهم نیست چون کامپایلر خودش ترتیب اجرا رو از رو ساختار گراف درمیاره
\end{dfnbox}

\begin{dfnbox}{بلوک config}{block-config}
بلوک \code{config} واسه تعیین یه سری مقادیر تو بلوک \code{\_\_main\_\_} کد تولیدشده به کار میره
مثلاً \code{batch\_size} پیش‌فرضش ۱ هست و \code{device} هم \code{'cpu'}
اگه کلید \code{loss} رو تنظیم کنیم، مولد کد برامون یه حلقه آموزش کامل با بهینه‌ساز و تابع خطا می‌سازه
\end{dfnbox}

\section{گرامر}

\begin{dfnbox}{آمار گرامر NNGraph.g4}{grammar-stats}
\begin{itemize}[noitemsep]
    \item فایل گرامر: \code{NNGraph.g4} (۵۶ خط)
    \item قوانین تجزیه‌گر: ۱۴ تا قانون (\LR{snake\_case})
    \item توکن‌های واژه‌نگار: ۲۰ تا توکن (\LR{ALL\_CAPS})
\end{itemize}
\end{dfnbox}

قوانین کامل گرامر هم ایناست:

\begin{codebox}{گرامر NNGraph.g4}{code-grammar}
\begin{amzcode}{antlr}
grammar NNGraph;
start     : model_block graph_block config_block? EOF;
model_block : MODEL ID '{' input_decl output_decl '}';
input_decl  : INPUT ID ':' TENSOR shape_expr;
output_decl : OUTPUT ID;
graph_block : GRAPH '{' (node_decl | edge_decl)* '}';
node_decl   : NODE ID ':' layer_expr;
layer_expr  : ID '(' param_list? ')';
param_list  : param (',' param)*;
param       : ID '=' value;
edge_decl   : EDGE ID ARROW ID ('[' LABEL '=' STRING ']')?;
config_block: CONFIG '{' config_entry* '}';
config_entry: ID '=' value;
value       : FLOAT_LITERAL | INT_LITERAL | BOOL_LITERAL
            | STRING | NONE | shape_expr;
shape_expr  : '(' INT_LITERAL (',' INT_LITERAL)* ')';
\end{amzcode}
\end{codebox}

\begin{notebox}
\textbf{نکته مهم طراحی:} واسه \code{input\_decl} ما از \code{TENSOR shape\_expr} استفاده کردیم نه \code{TENSOR '(' shape\_expr ')'} چون \code{shape\_expr} خودش پرانتز داره
اینطوری از یه باگ رایج که تو نسخه‌های اولیه‌مون داشتیم جلوگیری کردیم
\end{notebox}

\section{سیستم انواع}

\begin{center}
\begin{tabular}{lll}
\toprule
\textbf{نوع در DSL} & \textbf{نوع Python} & \textbf{مثال} \\
\midrule
\code{int}    & \code{int}        & \code{128} \\
\code{float}  & \code{float}      & \code{0.1} \\
\code{bool}   & \code{bool}       & \code{true} / \code{false} \\
\code{string} & \code{str}        & \code{"relu"} \\
\code{shape}  & \code{tuple[int]} & \code{(64, 32, 32)} \\
\code{None}   & \code{NoneType}   & \code{None} \\
\bottomrule
\end{tabular}
\end{center}

ما تمایز بین \code{int} و \code{float} رو همون زمان کامپایل چک می‌کنیم
مثلاً اگه بنویسیم \code{Dropout(p=1)} بهمون ارور می‌ده چون \code{p} منتظر یه \code{float} هست نه \code{int}

\section{انواع لایه}

\subsection{لایه‌های پارامتردار}

این لایه‌ها موقع ساخت \code{\_\_init\_\_} به شکل \code{self.<id> = nn.X(...)} تو کدمون جا می‌گیرن

\begin{center}
\begin{tabular}{llp{7.5cm}}
\toprule
\textbf{کلیدواژه DSL} & \textbf{کلاس PyTorch} & \textbf{پارامترهای اصلی} \\
\midrule
\code{Linear}        & \code{nn.Linear}             & \code{in\_features}, \code{out\_features}, \code{bias} \\
\code{Conv2d}        & \code{nn.Conv2d}             & \code{in\_ch}, \code{out\_ch}, \code{kernel}, \code{stride}, \code{padding} \\
\code{Conv1d}        & \code{nn.Conv1d}             & \code{in\_ch}, \code{out\_ch}, \code{kernel}, \code{stride} \\
\code{BatchNorm2d}   & \code{nn.BatchNorm2d}        & \code{num\_features} \\
\code{LayerNorm}     & \code{nn.LayerNorm}          & \code{normalized\_shape} \\
\code{MaxPool2d}     & \code{nn.MaxPool2d}          & \code{kernel}, \code{stride} \\
\code{AvgPool2d}     & \code{nn.AvgPool2d}          & \code{kernel}, \code{stride} \\
\code{Dropout}       & \code{nn.Dropout}            & \code{p} \\
\code{Flatten}       & \code{nn.Flatten}            & \code{start\_dim}, \code{end\_dim} \\
\code{Embedding}     & \code{nn.Embedding}          & \code{num\_embeddings}, \code{embedding\_dim} \\
\code{MultiHeadAttn} & \code{nn.MultiheadAttention} & \code{embed\_dim}, \code{num\_heads} \\
\code{LSTM}          & \code{nn.LSTM}               & \code{input\_size}, \code{hidden\_size}, \code{num\_layers} \\
\code{GRU}           & \code{nn.GRU}                & \code{input\_size}, \code{hidden\_size} \\
\bottomrule
\end{tabular}
\end{center}

\subsection{فعال‌سازی‌ها}

\code{ReLU}، \code{Sigmoid}، \code{Tanh}، \code{GELU}، \code{Softmax(dim)}، \code{LeakyReLU(negative\_slope)} و \code{ELU(alpha)}

\subsection{عملیات گراف (بدون ماژول)}

اینا دیگه تو \code{\_\_init\_\_} خطی تولید نمی‌کنن:

\begin{center}
\begin{tabular}{llp{7.5cm}}
\toprule
\textbf{کلیدواژه} & \textbf{پارامترها} & \textbf{کد تولیدشده} \\
\midrule
\code{Add}      & ---                          & \code{out = a + b} \\
\code{Concat}   & \code{dim: int}              & \code{out = torch.cat([a,b], dim=d)} \\
\code{Residual} & ---                          & \code{out = main + shortcut} \\
\code{Split}    & \code{chunks: int, dim: int} & \code{split\_0, split\_1, ... = torch.chunk(x,n,dim=d)} \\
\bottomrule
\end{tabular}
\end{center}

\subsection{پارامترهای سراسری}

یه پارامتر باحال به اسم \code{quant} داریم که رو هر گرهی می‌تونیم بذاریم
مقادیری که قبول می‌کنه ایناست: \code{"dynamic"}، \code{"static"} و \code{"qat"}

\begin{codebox}{نمونه استفاده از کوانتیزاسیون}{code-quant-example}
\begin{amzcode}{text}
node fc1 : Linear(in_features=784, out_features=128, quant="dynamic")
\end{amzcode}
\end{codebox}

\section{نحو یال}

\begin{dfnbox}{یال ساده}{edge-simple}
\begin{latin}
\texttt{edge src -> dst}
\end{latin}

اینطوری یه یال جهت‌دار از گره \code{src} به گره \code{dst} می‌سازیم
\end{dfnbox}

\begin{dfnbox}{یال برچسب‌دار}{edge-labeled}
\begin{latin}
\texttt{edge src -> dst [label="name"]}
\end{latin}

اینم واسه تعریف یه یال با برچسب رشته‌ایه
ما کلاً برچسب‌ها رو واسه دو تا کار استفاده می‌کنیم:
\begin{itemize}
    \item \code{"shortcut"}: واسه اینکه یال اتصال کوتاه به گره \code{Residual()} رو نشون بدیم
    \item \code{"residual\_*"}: واسه اتصال کوتاه به گرهی که دو تا ورودی داره (مثل الگوی \LR{transformer})
\end{itemize}
\end{dfnbox}

\section{مثال کامل: MLP}

فایل \code{.nng} واسه یه شبکه سه‌لایه این شکلی می‌شه:

\begin{codebox}{mlp.nng}{code-mlp-full}
\begin{amzcode}{text}
model MLP {
    input  x   : tensor(784)
    output out
}

graph {
    node fc1   : Linear(in_features=784, out_features=256)
    node relu1 : ReLU()
    node fc2   : Linear(in_features=256, out_features=128)
    node relu2 : ReLU()
    node drop  : Dropout(p=0.3)
    node fc3   : Linear(in_features=128, out_features=10)
    node out   : Softmax(dim=1)

    edge x     -> fc1
    edge fc1   -> relu1
    edge relu1 -> fc2
    edge fc2   -> relu2
    edge relu2 -> drop
    edge drop  -> fc3
    edge fc3   -> out
}

config {
    batch_size = 64
    device = "cuda"
}
\end{amzcode}
\end{codebox}

و اینم کد \LR{PyTorch} که برامون تولید می‌کنه:

\begin{codebox}{کد PyTorch تولیدشده برای MLP}{code-mlp-output}
\begin{amzcode}{python}
import torch
import torch.nn as nn


class MLP(nn.Module):
    def __init__(self):
        super(MLP, self).__init__()
        self.fc1   = nn.Linear(in_features=784, out_features=256)
        self.relu1 = nn.ReLU()
        self.fc2   = nn.Linear(in_features=256, out_features=128)
        self.relu2 = nn.ReLU()
        self.drop  = nn.Dropout(p=0.3)
        self.fc3   = nn.Linear(in_features=128, out_features=10)
        self.out   = nn.Softmax(dim=1)

    def forward(self, x):
        x = self.fc1(x)
        x = self.relu1(x)
        x = self.fc2(x)
        x = self.relu2(x)
        x = self.drop(x)
        x = self.fc3(x)
        x = self.out(x)
        return x


if __name__ == '__main__':
    device = torch.device('cuda')
    model  = MLP().to(device)
    x      = torch.randn(64, 784).to(device)
    print(model(x).shape)
\end{amzcode}
\end{codebox}

% ================================================================
\chapter{پیاده‌سازی}

\section{معماری کامپایلر}

\begin{dfnbox}{خط لوله کامپایل}{pipeline}
ما خط لوله کامپایلر رو تو شش تا مرحله چیدیم:
\begin{enumerate}
    \item \textbf{تجزیه}: \code{FileStream} $\to$ \code{NNGraphLexer} $\to$ \code{NNGraphParser.program()} $\to$ درخت تجزیه
    \item \textbf{تحلیل معنایی}: \code{ParseTreeWalker} + \code{NNGraphCustomListener} $\to$ مدل گراف
    \item \textbf{مرتب‌سازی توپولوژیک}: اینجا با الگوریتم \LR{Kahn} ترتیب اجرا رو درمیاریم
    \item \textbf{استنتاج شکل}: شکل تنسورها رو تو گراف پخش می‌کنیم و اگه مشکلی بود مچش رو می‌گیریم (اختیاری: \code{-{}-show-shapes})
    \item \textbf{تولید کد}: \code{CodeGenerator.generate()} $\to$ رشته پایتونی
    \item \textbf{خروجی}: تو این مرحله می‌ریزیم تو فایل \code{.py} یا یه خروجی \LR{Graphviz DOT} می‌گیریم یا کلاً خروجی \LR{ONNX} می‌دیم بیرون
\end{enumerate}
\end{dfnbox}

\subsection{مسئولیت فایل‌ها}

\begin{center}
\begin{tabular}{ll}
\toprule
\textbf{فایل} & \textbf{نقش} \\
\midrule
\code{NNGraph.g4}             & منبع حقیقت گرامر (۵۶ خط) \\
\code{gen/NNGraphLexer.py}    & تولیدشده توسط ANTLR4 \\
\code{gen/NNGraphParser.py}   & تولیدشده توسط ANTLR4 \\
\code{gen/NNGraphListener.py} & listener پایه --- نباید ویرایش شود \\
\code{custom\_listener.py}    & تحلیل معنایی \\
\code{code\_generator.py}     & تولید کد PyTorch \\
\code{shape\_inference.py}    & استنتاج شکل تنسور در کامپایل \\
\code{graph\_viz.py}          & تولید نمودار Graphviz DOT \\
\code{onnx\_export.py}        & صادرکردن مدل به فرمت ONNX \\
\code{main.py}                & رابط CLI و \code{compile\_nng()} \\
\bottomrule
\end{tabular}
\end{center}

\section{تحلیل‌گر معنایی}

\begin{dfnbox}{NNGraphCustomListener}{custom-listener}
ما یه کلاس \code{NNGraphCustomListener(NNGraphListener)} تو فایل \code{custom\_listener.py} داریم که کل بار تحلیل معنایی رو دوششه
وضعیت داخلیش هم اینطوریه:
\begin{itemize}[noitemsep]
    \item \code{self.nodes}: \code{dict[id, \{type, params, line\}]}
    \item \code{self.edges}: \code{list[\{src, dst, label, line\}]}
    \item \code{self.config}: \code{dict[key, value]}
\end{itemize}
\end{dfnbox}

\subsection{انتشار مقدار در درخت}

متدهای \code{exitShape\_expr} و \code{exitValue} مقادیر رو روی خود \LR{context} می‌نویسن:

\begin{codebox}{انتشار مقدار ctx.pvalue در exitShape\_expr}{code-propagation}
\begin{amzcode}{python}
def exitShape_expr(self, ctx):
    ints = [int(ctx.INT_LITERAL(i).getText())
            for i in range(ctx.INT_LITERAL().__len__())]
    ctx.pvalue = tuple(ints)
    ctx.ptype  = tuple
\end{amzcode}
\end{codebox}

این روش به ما اطمینان می‌ده که وقتی از یه \LR{context} میایم بیرون، مقادیر بچه‌هاش دم دستمون باشن

\subsection{قوانین اعتبارسنجی}

ما هشت تا قانون اعتبارسنجی داریم که تو دو تا فاز اجراشون می‌کنیم:

\textbf{فاز اول (همون موقعی که داریم درخت رو بالا پایین می‌کنیم):}
\begin{enumerate}
    \item \textbf{شناسه تکراری گره}: تو \code{exitNode\_decl} چکش می‌کنیم
    \item \textbf{ارجاع نامعتبر در یال}: تو \code{exitEdge\_decl}
    \item \textbf{نوع لایه ناشناس}: تو \code{exitNode\_decl}
    \item \textbf{عدم تطابق نوع پارامتر}: بازم تو \code{exitNode\_decl}
\end{enumerate}

\textbf{فاز دوم (وقتی گرافمون کامل شد تو \code{exitGraph\_block}):}
\begin{enumerate}
    \setcounter{enumi}{4}
    \item \textbf{گره خروجی اعلان‌نشده}: تو \code{\_validate\_output\_declared}
    \item \textbf{گره یتیم}: تو \code{\_validate\_orphans} مچش رو می‌گیریم
    \item \textbf{آریتی residual != 2}: تو \code{\_validate\_residual\_arity}
    \item \textbf{تشخیص دور با kahn}: تو \code{\_validate\_no\_cycles}
    \item \textbf{قابلیت دسترسی BFS}: تو \code{\_validate\_reachability}
\end{enumerate}

\subsection{قالب پیام خطا}

\begin{codebox}{قالب خطا}{code-error-format}
\begin{amzcode}{text}
[line L:C] SemanticError: <message>
[line L] ShapeError at '<node>': <message>
\end{amzcode}
\end{codebox}

اینجا \code{L} همون شماره خطه و \code{C} هم ستون توکنمونه
اگه به هر کدوم از این خطاها بخوریم، کامپایل متوقف می‌شه و با کد خروج ۱ می‌پریم بیرون

\subsection{\_FakeCtx برای خطاهای پس از پیمایش}

خطاهایی که تو \code{exitGraph\_block} پیش میان \LR{context} زنده ندارن
واسه همین ما یه کلاس جمع‌وجور به اسم \code{\_FakeCtx} ساختیم که شماره خط رو از گره برداره و با خودش بیاره:

\begin{codebox}{کلاس \_FakeCtx برای حفظ شماره خط}{code-fakectx}
\begin{amzcode}{python}
class _FakeCtx:
    class start:
        line   = info["line"]
        column = 0
\end{amzcode}
\end{codebox}

\subsection{جریان اجرای شنونده}

متد \code{ParseTreeWalker.walk(listener, parse\_tree)} یه پیمایش عمق‌اول روی درختمون می‌زنه و قبل از اینکه بره سراغ بچه‌ها متد \code{enterX(ctx)} رو صدا می‌زنه و بعدش \code{exitX(ctx)} رو

\begin{dfnbox}{فراخوانی واکر در main.py}{walker-invoke}
\begin{amzcode}{python}
listener = NNGraphCustomListener()
listener.rule_names = parser.ruleNames
walker = ParseTreeWalker()
walker.walk(listener, parse_tree)
\end{amzcode}
\end{dfnbox}

اینجا شنونده سفارشی ما فقط متدهای \code{exit*} رو بازنویسی می‌کنه
متدهای \code{enter*} رو هم همونجوری بدون تغییر از کلاس پایه \code{NNGraphListener} نگه می‌داریم

\subsubsection{انتشار ویژگی‌ها (پایین به بالا)}

ما مقادیر رو از طریق ویژگی‌های کانتکست می‌فرستیم بالا تو درخت:

\begin{codebox}{متد exitValue مقدار پارس‌شده را در ctx.pvalue ذخیره می‌کند}{code-exitvalue}
\begin{amzcode}{python}
def exitValue(self, ctx):
    if ctx.FLOAT_LITERAL():
        ctx.pvalue = float(ctx.FLOAT_LITERAL().getText())
        ctx.ptype  = float
    elif ctx.INT_LITERAL():
        ctx.pvalue = int(ctx.INT_LITERAL().getText())
        ctx.ptype  = int
    elif ctx.BOOL_LITERAL():
        ctx.pvalue = ctx.BOOL_LITERAL().getText() == 'true'
        ctx.ptype  = bool
    # ... other cases
\end{amzcode}
\end{codebox}

اینطوری قوانین والد هم می‌تونن ویژگی‌های بچه‌هاشون رو بخونن:

\begin{codebox}{متد exitParam مقدار ctx.value().pvalue را از فرزند می‌خواند}{code-exitparam}
\begin{amzcode}{python}
def exitParam(self, ctx):
    ctx.param_name  = ctx.ID().getText()
    ctx.param_value = ctx.value().pvalue
    ctx.param_type  = ctx.value().ptype
\end{amzcode}
\end{codebox}

\subsubsection{استخراج معنایی و شرح تفصیلی رفتار متدها}

شنونده سفارشی ما با یه سری متد \code{exit*} کار می‌کنه که از پایین به بالا میان
اینجا براتون باز می‌کنم که هر تابع شنونده دقیقاً چه کاری می‌کنه:

\begin{itemize}
    \item \textbf{\code{exitShape\_expr(ctx)}}: این میاد لیستی از ابعاد رو از توکن‌های \code{INT\_LITERAL} پارس می‌کنه و یه \code{tuple} می‌سازه، بعدم می‌ریزه تو \code{ctx.pvalue} و نوعشم می‌ذاره \code{tuple}
    \item \textbf{\code{exitValue(ctx)}}: مقادیر لیتراال رو ارزیابی می‌کنه و تبدیلشون می‌کنه به معادل پایتونی‌شون و می‌ریزه تو \code{ctx.pvalue}
    \item \textbf{\code{exitModel\_block(ctx)}}: شناسه مدل شبکه رو از \code{ctx.ID().getText()} درمیاره و می‌ذاره تو \code{self.model\_name}
    \item \textbf{\code{exitInput\_decl(ctx)}}: شناسه متغیر ورودی رو تو \code{self.input\_id} و شکل پارس‌شده‌ش رو تو \code{self.input\_shape} ذخیره می‌کنه
    تازه یه شبه‌گره \code{"\_\_input\_\_"} هم می‌سازه واسه اعتبارسنجی
    \item \textbf{\code{exitOutput\_decl(ctx)}}: شناسه خروجی رو نگه می‌داره تا بدونیم ته‌مون کجاست
    \item \textbf{\code{exitParam(ctx)}}: پارامترها رو هندل می‌کنه و شناسه و مقدار و نوعشون رو ذخیره می‌کنه
    \item \textbf{\code{exitLayer\_expr(ctx)}}: پیکربندی‌های سازنده لایه رو می‌گیره و پارامترهاشون رو مپ می‌کنه
    \item \textbf{\code{exitNode\_decl(ctx)}}: گره‌ها رو ثبت می‌کنه و چک می‌کنه شناسه تکراری نباشه و پارامترها هم به \code{LAYER\_SCHEMA} بخورن
    \item \textbf{\code{exitEdge\_decl(ctx)}}: یال‌ها رو پردازش می‌کنه و چک می‌کنه که مبدأ و مقصد تو \code{self.nodes} باشن و بعد یال رو ثبت می‌کنه
    \item \textbf{\code{exitConfig\_entry(ctx)}}: تنظیمات سراسری مثل \code{device} یا \code{batch\_size} رو می‌گیره می‌ذاره تو \code{self.config}
    \item \textbf{\code{exitGraph\_block(ctx)}}: وقتی گراف کامل شد می‌پره وسط و اون پنج تا مرحله اعتبارسنجی که قبلاً گفتم رو اجرا می‌کنه
    \item \textbf{\code{exitEveryRule(ctx)}}: خروج از قوانین رو ردگیری می‌کنه و اگه لازم بود \code{make\_ast\_subtree} رو صدا می‌زنه
    \item \textbf{\code{exitStart(ctx)}}: کار درخت \LR{AST} رو نهایی می‌کنه
\end{itemize}

\subsubsection{الگوریتم‌های اعتبارسنجی گراف}

متد \code{exitGraph\_block} وقتی گراف ساخته شد این پنج تا کار رو می‌کنه:

\textbf{۱. تشخیص گره یتیم} (\code{\_validate\_orphans}):

\begin{codebox}{الگوریتم تشخیص گره یتیم}{code-orphan}
\begin{amzcode}{python}
referenced = set()
for e in self.edges:
    referenced.add(e["src"])
    referenced.add(e["dst"])
for node_id, info in self.nodes.items():
    if info["type"] == "__input__":
        continue
    if node_id not in referenced:
        error(f"Node '{node_id}' is declared but never referenced.")
\end{amzcode}
\end{codebox}

اینجا شناسه همه گره‌هایی که تو یال‌ها هستن رو جمع می‌کنیم
هر گرهی که تعریف شده ولی تو این یال‌ها نیست، یتیم حساب می‌شه

\textbf{۲. تشخیص دور} (\code{\_validate\_no\_cycles}):

اینجا رفتیم سراغ الگوریتم \LR{Kahn}
اگه بعد از پردازش، گرهی بمونه که \code{in\_degree > 0} باشه، یعنی یه جایی دور داریم

\begin{codebox}{الگوریتم Kahn برای تشخیص دور}{code-kahn}
\begin{amzcode}{python}
in_degree = dict.fromkeys(self.nodes, 0)
adj       = {n: [] for n in self.nodes}
for e in self.edges:
    adj[e["src"]].append(e["dst"])
    in_degree[e["dst"]] += 1

queue   = deque(n for n, d in in_degree.items() if d == 0)
visited = 0
while queue:
    node = queue.popleft()
    visited += 1
    for nxt in adj[node]:
        in_degree[nxt] -= 1
        if in_degree[nxt] == 0:
            queue.append(nxt)

if visited != len(self.nodes):
    remaining = [n for n, d in in_degree.items() if d > 0]
    error(f"Cycle detected involving nodes: {remaining}.")
\end{amzcode}
\end{codebox}

\textbf{۳. قابلیت دسترسی} (\code{\_validate\_reachability}):

یه جستجوی \LR{BFS} از گره ورودی می‌زنیم و انتظار داریم همه گره‌ها قابل دسترسی باشن

\begin{codebox}{بررسی قابلیت دسترسی با BFS}{code-bfs}
\begin{amzcode}{python}
adj = {n: [] for n in self.nodes}
for e in self.edges:
    adj[e["src"]].append(e["dst"])

visited = set()
queue   = deque([self.input_id])
while queue:
    node = queue.popleft()
    if node in visited:
        continue
    visited.add(node)
    for nxt in adj[node]:
        queue.append(nxt)

for node_id, info in self.nodes.items():
    if node_id not in visited and info["type"] != "__input__":
        error(f"Node '{node_id}' is not reachable from input.")

if self.output_id not in visited:
    error(f"Output node '{self.output_id}' is not reachable.")
\end{amzcode}
\end{codebox}

اینطوری بخش‌های جدا افتاده رو پیدا می‌کنیم

\begin{notebox}
دقت کنین که تشخیص دور و قابلیت دسترسی دو تا چیز کاملاً جدان
ممکنه گرافمون دور نداشته باشه ولی یه تیکه‌ش کلاً جدا باشه، یا همه جاش وصل باشه ولی دور داشته باشه
ما واسه \LR{DAG} شبکه‌مون به جفتش نیاز داریم
\end{notebox}

\section{مولد کد}

\begin{dfnbox}{CodeGenerator}{codegen}
ما کلاس \code{CodeGenerator} رو تو \code{code\_generator.py} تو چهار تا گام زدیم:
\begin{enumerate}[noitemsep]
    \item \code{\_compute\_graph()}: ساخت لیست‌های مجاورت و مرتب‌سازی توپولوژیک
    \item \code{\_assign\_vars()}: تخصیص نام متغیر پایتونی به هر گره
    \item \code{\_emit\_init()} + \code{\_emit\_forward()}: تولید خطوط کد
    \item \code{\_assemble()}: هم‌زدن همه اینا با هم تو یه سورس نهایی پایتون
\end{enumerate}
\end{dfnbox}

\section{الگوریتم تخصیص متغیر}

\begin{dfnbox}{اصل اساسی تخصیص متغیر}{var-invariant}
ما \code{var\_at[n] = 'x'} رو فقط موقعی می‌ذاریم که \code{out\_degree[predecessor] == 1} باشه

وقتی یه گره قراره به چند جا بره (\code{out\_degree > 1})، یعنی شاخه‌های دیگه هنوز بهش نیاز دارن
اگه همونجا \code{x} رو بازنویسی کنیم که کارمون زاره
واسه همین، گره بعدی همون اسم شناسه خودش رو به عنوان متغیر می‌گیره
\end{dfnbox}

ما تو \code{\_assign\_vars} چهار تا قانون کلی داریم:

\begin{codebox}{قوانین تخصیص متغیر در \_assign\_vars}{code-var-rules}
\begin{amzcode}{text}
Rule 1 - Input node:
    var_at[input_id] = 'x'

Rule 2 - Merge node (in_degree > 1):
    var_at[n] = 'x'       if out_degree[n] <= 1
    var_at[n] = node_id   if out_degree[n] > 1

Rule 3 - Branch start (predecessor.out_degree > 1):
    var_at[n] = node_id

Rule 4 - Linear continuation (predecessor.out_degree == 1):
    var_at[n] = var_at[predecessor]
\end{amzcode}
\end{codebox}

\subsection{ردیابی MLP (زنجیره خطی)}

\begin{exbox}{ردیابی متغیر برای MLP}{trace-mlp}
توپولوژی شبکه این شکلیه:
\begin{latin}
\texttt{x -> fc1 -> relu1 -> fc2 -> relu2 -> drop -> fc3 -> out}
\end{latin}

اینجا هر گره قبلی \code{out\_degree == 1} داره، پس طبق قانون ۴ همیشه \code{var\_at = 'x'} می‌شه
و تو \code{forward} همه‌جا از همون \code{x} استفاده می‌کنیم
\end{exbox}

\subsection{ردیابی ResBlock (فن‌اوت سپس ادغام)}

\begin{exbox}{ردیابی متغیر برای ResBlock}{trace-resblock}
ما اینجا \code{x} رو هم به \code{conv1} و هم به \code{skip} می‌فرستیم، پس \code{out\_degree[x] = 2} می‌شه

\begin{itemize}
    \item واسه \code{conv1}: قانون ۳ میگه \code{var\_at = 'conv1'}
    \item واسه بقیه: قانون ۴ میگه همه همون \code{'conv1'} رو می‌گیرن
    \item واسه \code{skip}: قانون ۲ میگه \code{var\_at = 'x'}
\end{itemize}

اینم کدی که تولید می‌کنه:
\begin{codebox}{forward pass تولیدشده برای ResBlock}{code-resblock-forward}
\begin{amzcode}{python}
conv1 = self.conv1(x)
conv1 = self.bn1(conv1)
conv1 = self.relu1(conv1)
conv1 = self.conv2(conv1)
conv1 = self.bn2(conv1)
x     = conv1 + x       # Residual node
x     = self.flatten(x)
\end{amzcode}
\end{codebox}
\end{exbox}

\subsection{ردیابی InceptionBlock (دو شاخه موازی)}

\begin{exbox}{ردیابی متغیر برای InceptionBlock}{trace-inception}
اینجا هم \code{out\_degree[x] = 2} هست

\begin{itemize}
    \item واسه \code{convA1} و \code{convB1}: قانون ۳ بهمون \code{'convA1'} و \code{'convB1'} می‌ده
    \item واسه \code{convA2} و \code{convB2}: قانون ۴ رو داریم
    \item واسه \code{cat1}: قانون ۲ دوباره بهمون \code{'x'} می‌ده
\end{itemize}

کدش هم اینطوری درمیاد:
\begin{codebox}{forward pass تولیدشده برای InceptionBlock}{code-inception-forward}
\begin{amzcode}{python}
convA1 = self.convA1(x)
convA1 = self.convA2(convA1)
convB1 = self.convB1(x)
convB1 = self.convB2(convB1)
x      = torch.cat([convA1, convB1], dim=1)  # Concat node
x      = self.bn(x)
x      = self.out(x)
\end{amzcode}
\end{codebox}
\end{exbox}

\subsection{ردیابی TransformerEncoder (دو residual ضمنی)}

\begin{exbox}{ردیابی متغیر برای TransformerEncoder}{trace-transformer}
اینجا وضعیت جالب‌تره، \code{out\_degree[x] = 2} و \code{out\_degree[norm1] = 2}

\begin{itemize}
    \item واسه \code{attn}: قانون ۳
    \item واسه \code{drop1}: قانون ۴
    \item واسه \code{norm1}: قانون ۲ بهمون \code{'norm1'} می‌ده
    \item واسه \code{ff1}: قانون ۳
    \item واسه \code{gelu1} و \code{drop2} و \code{ff2}: قانون ۴
    \item واسه \code{norm2}: قانون ۲ بهمون \code{'x'} می‌ده
\end{itemize}

خروجیش این شکلیه:
\begin{codebox}{forward pass تولیدشده برای TransformerEncoder}{code-transformer-forward}
\begin{amzcode}{python}
attn, _ = self.attn(x, x, x)    # MultiHeadAttn triple-pass
attn     = self.drop1(attn)
norm1    = self.norm1(x + attn)  # implicit residual_1
ff1      = self.ff1(norm1)
ff1      = self.gelu1(ff1)
ff1      = self.drop2(ff1)
ff1      = self.ff2(ff1)
x        = self.norm2(norm1 + ff1) # implicit residual_2
x        = self.out(x)
x        = self.flatten(x)
return x
\end{amzcode}
\end{codebox}
\end{exbox}

\subsection{موارد خاص تولید کد}

\begin{center}
\begin{tabular}{llp{7.5cm}}
\toprule
\textbf{الگو} & \textbf{شرط} & \textbf{کد تولیدشده} \\
\midrule
\LR{MultiHeadAttn}  & نوع لایه     & \code{out, \_ = self.attn(x, x, x)} \\
\LR{LSTM / GRU}     & نوع لایه     & \code{out, \_ = self.lstm(x)} \\
\LR{Residual}       & بدون ماژول   & \code{out = main + shortcut} \\
\LR{Add}            & بدون ماژول   & \code{out = a + b} \\
\LR{Concat}         & بدون ماژول   & \code{out = torch.cat([a,b], dim=d)} \\
\LR{Split}          & بدون ماژول   & \code{split\_0, split\_1, ... = torch.chunk(x,n,dim=d)} \\
\LR{Residual ضمنی}  & \code{residual\_*} label & \code{out = self.norm(skip + main)} \\
\bottomrule
\end{tabular}
\end{center}

\begin{notebox}
\textbf{نگاشت نام پارامتر:} ما تو DSL اسما رو کوتاه کردیم ولی تو تولید کد به معادل‌های پایتونیشون برش گردوندیم
مثلاً \code{in\_ch} می‌شه \code{in\_channels}
\end{notebox}

\subsection{ساختار کد خروجی}

\begin{codebox}{الگوی کد خروجی تولیدشده}{code-output-template}
\begin{amzcode}{python}
import torch
import torch.nn as nn


class <ModelName>(nn.Module):
    def __init__(self):
        super(<ModelName>, self).__init__()
        # one line per parameterized layer node

    def forward(self, x):
        # one line per node in topological order
        return <output_var>


if __name__ == '__main__':
    device = torch.device('<device>')
    model  = <ModelName>().to(device)
    x      = torch.randn(<batch_size>, <input_shape>).to(device)
    print(model(x).shape)
\end{amzcode}
\end{codebox}

\section{استنتاج شکل}

\begin{dfnbox}{shape inference}{shape-inference}
ما ماژول \code{shape\_inference.py} رو نوشتیم که شکل تنسور رو از ورودی تو کل گراف پخش کنه
قوانینش هم سادس، مثلاً \LR{Linear} بعد آخر رو عوض می‌کنه یا \LR{Flatten} ابعاد رو با هم یکی می‌کنه
\end{dfnbox}

اگه ابعاد به هم نخورن، خیلی شیک خطای \code{ShapeError} می‌دیم و کامپایل رو قطع می‌کنیم
با پرچم \code{-{}-show-shapes} هم می‌شه شکل‌ها رو دید

\section{تجسم Graphviz}

تو ماژول \code{graph\_viz.py} ما خروجی \LR{Graphviz DOT} می‌دیم بیرون
به نظرم رنگ‌بندیشم باحاله، مثلاً ورودی سبزه و لایه‌های نرمال‌سازی بنفش
گره خروجی رم با دو تا خط مشخص کردیم و یال‌های برچسب‌دار رو خط‌چین زدیم

\begin{codebox}{تولید و رندر نمودار}{code-graphviz}
\begin{amzcode}{bash}
python main.py -i input/inception.nng --graph-viz output/inception.dot
dot -Tpng output/inception.dot -o output/inception.png
\end{amzcode}
\end{codebox}

\section{داربست آموزش}

اگه \code{loss} رو تو تنظیمات بدیم، کد ژنراتورمون یه حلقه آموزش مشتی برامون سرهم می‌کنه
توابع خطای معروف مثل \code{CrossEntropyLoss} و \code{MSELoss} و بهینه‌سازایی مثل \code{Adam} و \code{SGD} رو هم پشتیبانی می‌کنیم

\section{حاشیه‌نویسی کوانتیزاسیون}

پارامتر سراسری \code{quant} رو همه‌جا می‌شه استفاده کرد
سه تا حالت \code{"dynamic"} و \code{"static"} و \code{"qat"} رو ساپورت می‌کنیم
اگه کوانتیزاسیون فعال باشه، اون \code{QuantStub} و دوستاشم به کدمون اضافه می‌شن

\section{صادرکردن ONNX}

ماژول \code{onnx\_export.py} مدلمون رو نمونه‌سازی می‌کنه و یه خروجی تمیز \LR{ONNX} می‌ده که با محور دسته‌ای پویاست

\begin{codebox}{صادرکردن مدل به ONNX}{code-onnx}
\begin{amzcode}{bash}
python main.py -i input/mlp.nng --onnx output/mlp.onnx
\end{amzcode}
\end{codebox}

% ================================================================
\chapter{ارزیابی}

\section{مجموعه آزمایش}

\begin{dfnbox}{آمار آزمایش}{test-stats}
\begin{itemize}[noitemsep]
    \item تعداد کل تست‌ها: \textbf{۵۵}
    \item فریم‌ورک: \code{pytest} نسخه ۸.۴.۲
    \item زمان اجرا: حدود ۱.۰۹ ثانیه
\end{itemize}
\end{dfnbox}

\begin{center}
\begin{tabular}{lcp{6cm}}
\toprule
\textbf{فایل} & \textbf{تست‌ها} & \textbf{پوشش} \\
\midrule
\code{test\_parser.py}   & ۸  & گرامر بدون خطای نحوی برای ۴ نمونه؛ config، برچسب، شکل \\
\code{test\_listener.py} & ۱۴ & ۴ ورودی معتبر؛ ۱۰ شرط خطا با \code{SystemExit} \\
\code{test\_codegen.py}  & ۱۹ & نام کلاس‌های nn، مقادیر پارامتر، الگوهای forward \\
\code{test\_e2e.py}      & ۱۴ & کامپایل کامل + ایجاد نمونه + forward pass با شکل صحیح \\
\bottomrule
\end{tabular}
\end{center}

اینم دستورش:

\begin{codebox}{اجرای مجموعه آزمایش}{code-run-tests}
\begin{amzcode}{bash}
python -m pytest tests/ -v
\end{amzcode}
\end{codebox}

\section{نتایج forward pass}

\begin{exbox}{MLP}{eval-mlp}
\begin{itemize}[noitemsep]
    \item ورودی: \code{tensor(784)} --- خروجی: \code{(batch, 10)}
    \item آزمایش: \code{softmax} ردیف‌ها به ۱ جمع می‌شن
\end{itemize}
\end{exbox}

\begin{exbox}{ResBlock}{eval-resblock}
\begin{itemize}[noitemsep]
    \item ورودی: \code{tensor(64, 32, 32)} --- خروجی: \code{(batch, 65536)}
    \item آزمایش: شکلش درسته و تو حالت \code{eval} خروجی قطعیه
\end{itemize}
\end{exbox}

\begin{exbox}{InceptionBlock}{eval-inception}
\begin{itemize}[noitemsep]
    \item ورودی: \code{tensor(32, 56, 56)} --- خروجی: \code{(batch, 48, 56, 56)}
    \item شاخه A با ۱۶ و شاخه B با ۳۲ کانال ادغام می‌شن و ۴۸ کانال می‌دن
\end{itemize}
\end{exbox}

\begin{exbox}{TransformerEncoder}{eval-transformer}
\begin{itemize}[noitemsep]
    \item ورودی: \code{tensor(512, 128)} --- خروجی: \code{(1024, 128)}
    \item دو تا زیرلایه residual ضمنی داریم
\end{itemize}
\end{exbox}

\section{اثبات درستی متغیرها}

\begin{dfnbox}{اثبات عدم بازنویسی تنسور}{var-correctness}
ما واسه هر نمونه، \code{data\_ptr()} تنسور ورودی رو چک کردیم:
\begin{itemize}[noitemsep]
    \item \textbf{ResBlock}: اینجا \code{x.data\_ptr()} دست‌نخورده مونده
    \item \textbf{Inception}: تو جفت شاخه‌ها \code{x.data\_ptr()} یکسانه
    \item \textbf{Transformer}: تو تمام مسیرهای residual هم مقادیر تغییر نکردن
\end{itemize}

این دقیقاً همون چیزیه که ما از قانونمون انتظار داشتیم
\end{dfnbox}

\section{آمار کدبیس}

\begin{center}
\begin{tabular}{lr}
\toprule
\textbf{فایل/دایرکتوری} & \textbf{خطوط} \\
\midrule
\code{NNGraph.g4}                    & ۵۶ \\
\code{custom\_listener.py}           & ۳۰۳ \\
\code{code\_generator.py}            & ۳۲۷ \\
\code{shape\_inference.py}           & ۱۹۳ \\
\code{graph\_viz.py}                 & ۹۰ \\
\code{onnx\_export.py}              & ۴۷ \\
\code{main.py}                       & ۱۲۱ \\
\code{tests/} (۴ فایل)              & ۵۷۵ \\
\code{required\_code\_collection/}   & $\approx$۱۰۰ \\
\midrule
\textbf{جمع}                         & $\approx$۱۸۱۲ \\
\bottomrule
\end{tabular}
\end{center}

\begin{itemize}[noitemsep]
    \item زبان: \LR{Python 3.14.3}
    \item \LR{ANTLR: 4.13.2}
    \item \LR{PyTorch: 2.10.0}
    \item \LR{pytest: 8.4.2}
\end{itemize}

% ================================================================
\chapter{تصمیمات طراحی}

\section{Listener در برابر Visitor}

الگوی \LR{Listener} خیلی راحت بهمون اجازه داد وضعیت رو نگه داریم و مستقیم رو \code{self.nodes} و \code{self.edges} کار کنیم

\section{load\_from\_listener() در برابر generate(traversal)}

واسه کامپایلر گراف، ساختار درختی و گراف خودش خیلی طبیعی‌تر از یه لیست خطی توکنه، واسه همین \code{load\_from\_listener} رو ترجیح دادیم

\section{نام‌گذاری متغیر با node\_id}

با اینکه می‌شد از اسمای تک‌حرفی استفاده کرد، ما رفتیم سراغ شناسه واقعی گره
چون هم قطعیه، هم خیلی خواناتر می‌شه و تو گرافای بزرگ هم کارمون رو راحت می‌کنه

% ================================================================
\chapter{نتیجه‌گیری و کارهای آینده}

\section{نتیجه‌گیری}

در کل ما یه \LR{NNGraph DSL} مشتی با شش مرحله کامپایل ساختیم
یه کامپایلر با ۱۴ قانون تجزیه‌گر، ۲۴ نوع لایه و کلی قوانین اعتبارسنجی داریم که حتی پیچیده‌ترین شاخه‌ها رو هم بدون باگ هندل می‌کنه
تستاشم رو ۴ تا معماری کامل ران کردیم و جواب داد

\section{قابلیت‌های پیاده‌سازی‌شده}

\begin{center}
\begin{tabular}{lp{10cm}}
\toprule
\textbf{قابلیت} & \textbf{توضیح} \\
\midrule
استنتاج شکل        & انتشار شکل تنسور و تشخیص عدم تطابق در کامپایل (\code{-{}-show-shapes}) \\
تجسم Graphviz       & خروجی نمودار گراف مدل به فرمت DOT (\code{-{}-graph-viz}) \\
داربست آموزش        & تولید حلقه آموزش با بهینه‌ساز و تابع خطا از بلوک \code{config} \\
کوانتیزاسیون        & متادیتا \code{quant} برای تبدیل مدل به نسخه‌های کوانتیزه \\
صادرکردن ONNX       & صادرکردن مدل به فرمت ONNX (\code{-{}-onnx}) \\
\bottomrule
\end{tabular}
\end{center}

\section{کارهای آینده}

\begin{center}
\begin{tabular}{lp{10cm}}
\toprule
\textbf{قابلیت} & \textbf{توضیح} \\
\midrule
زیرگراف‌های نام‌دار  & بلوک‌های ماکرو قابل استفاده مجدد \\
کنترل جریان پویا   & شاخه‌های \code{if}/\code{else} درون \code{forward} \\
\bottomrule
\end{tabular}
\end{center}

% ================================================================
\chapter*{منابع}
\markboth{منابع}{منابع}
\addcontentsline{toc}{chapter}{منابع}

\begin{enumerate}
    \item \LR{Parr, T. (2013). \textit{The Definitive ANTLR 4 Reference}. Pragmatic Bookshelf.}
    \item \LR{Paszke, A., et al. (2019). PyTorch: An Imperative Style, High-Performance Deep Learning Library. \textit{NeurIPS 2019}.}
    \item \LR{NADER Framework: Neural Architecture Design via Representation (referenced in \texttt{NNGraph\_DSL\_Specification.md}).}
    \item \LR{Zhang, A. M. \textit{amznotes LaTeX template}. \texttt{github.com/alexmingzhang/latex-notes-template}}
\end{enumerate}

\end{document}
